% !PRINTED none
% Abschnitt-opening page: no folio number visible in the running head on this page
% (chapter-opening layout). Neunter Abschnitt / Diskriminante / § 84 begins here.

\begin{center}
\gesperrt{Neunter Abschnitt.}

\textbf{\Large Diskriminante.}

\rule{0.3\linewidth}{0.4pt}
\end{center}

\begin{center}
\textbf{§ 84. \gesperrt{Definition der Diskriminanten.}}
\end{center}

In der Theorie der quadratischen Körper treten als Diskriminanten gewisse ganze rationale (positive oder negative) Zahlen auf, die, wenn sie gerade sind, durch 4 teilbar sind, und wenn sie ungerade sind, bei der Teilung durch 4 den Rest 1 lassen.

\gesperrt{Kronecker} hat solche Zahlen \glqq \gesperrt{Zahlen von Diskriminantenform}\grqq{} genannt. Wir wollen sie hier kurz Diskriminanten nennen, und uns nicht daran stoßen, daß dieses Wort in der Algebra noch mannigfache andere Bedeutungen hat. Wir definieren daher:

1. \gesperrt{Eine positive oder negative, von Null verschiedene ganze Zahl, die einer der beiden Kongruenzen}
\weq{1}{D\equiv 0,\quad D\equiv 1 \quad (\mathrm{mod}\ 4)}
\gesperrt{genügt, heißt eine Diskriminante.}

Jede Quadratzahl ist nach dieser Definition eine Diskriminante. Da diese aber eine gewisse Ausnahmestellung einnehmen, wollen wir sie, wenn die Unterscheidung nötig ist, \gesperrt{uneigentliche Diskriminanten nennen.}

2. \gesperrt{Das Produkt zweier und folglich beliebig vieler Diskriminanten ist wieder eine Diskriminante.}

3. \gesperrt{Eine Diskriminante, die (außer 1) keinen quadratischen Teiler enthält, nach dessen Absonderung eine Diskriminante übrig bleibt, heißt Stammdiskriminante.}

Stammdiskriminanten sollen in der Folge zum Unterschied von anderen mit $\Delta$ bezeichnet sein. Ist $D$ keine Stammdiskriminante, so gibt es eine und nur eine Quadratzahl $Q^2$ und eine Stammdiskriminante $\Delta$, so daß
\weq{2}{D = \Delta Q^2}
ist; $\Delta$ heißt dann der \gesperrt{Stamm von} $D$.
